\chapter{Standard logistic regression analysis}

In this section, we analyze the performance of a standard (non-prior-weighted) Logistic Regression model on the project dataset. The goal is to evaluate the impact of the regularization parameter $\lambda$ on both the discriminative capability (minDCF) and the calibration (actDCF) of the model for a target application characterized by an effective prior $\tilde{\pi} = \pi_T = 0.1$.

\section{Experimental results}

The model was trained using 13 values of $\lambda$ logarithmically spaced between $10^{-4}$ and $10^{2}$. For each value, we computed the Minimum DCF (minDCF) and Actual DCF (actDCF). To evaluate the actual Bayes decisions, we converted the raw logistic regression scores into valid Log-Likelihood Ratios (LLRs) by subtracting the empirical prior log-odds of the training set.

The numerical results are summarized in Table \ref{tab:std_lr_results} and the visual trend is plotted in Figure \ref{fig:lr_standard_plot}.

\begin{table}[h]
    \centering
    \caption{Standard Logistic Regression Performance ($\pi_T = 0.1$)}
    \label{tab:std_lr_results}
    \begin{tabular}{@{}ccc@{}}
        \toprule
        \textbf{Regularization Coefficient ($\lambda$)} & \textbf{minDCF} & \textbf{actDCF} \\ \midrule
        $1.0000 \times 10^{-4}$ & 0.3640 & 0.4021 \\
        $3.1623 \times 10^{-4}$ & 0.3650 & 0.4051 \\
        $1.0000 \times 10^{-3}$ & 0.3650 & 0.4130 \\
        $3.1623 \times 10^{-3}$ & 0.3641 & 0.4297 \\
        $1.0000 \times 10^{-2}$ & \textbf{0.3611} & 0.4568 \\
        $3.1623 \times 10^{-2}$ & 0.3621 & 0.5805 \\
        $1.0000 \times 10^{-1}$ & 0.3641 & 0.8522 \\
        $3.1623 \times 10^{-1}$ & 0.3640 & 0.9950 \\
        $1.0000 \times 10^{0}$  & 0.3640 & 1.0000 \\
        $3.1623 \times 10^{0}$  & 0.3640 & 1.0000 \\
        $1.0000 \times 10^{1}$  & 0.3640 & 1.0000 \\
        $3.1623 \times 10^{1}$  & 0.3640 & 1.0000 \\
        $1.0000 \times 10^{2}$  & 0.3640 & 1.0000 \\
        \bottomrule
    \end{tabular}
\end{table}

\begin{figure}[h]
    \centering
    \includegraphics[width=0.85\textwidth]{LR_Standard_DCF_vs_Lambda.png}
    \caption{Actual DCF and Minimum DCF as a function of the regularization parameter $\lambda$ for the Standard Logistic Regression model ($\pi_T = 0.1$).}
    \label{fig:lr_standard_plot}
\end{figure}

The optimal discriminative capability is reached at $\lambda = 10^{-2}$, yielding a minDCF of \textbf{0.3611}. However, at this specific operational point, the actual DCF is noticeably higher at $0.4568$, indicating a baseline calibration gap.

\section{Analysis and considerations}

\subsection*{Does the regularization coefficient affect the two metrics differently?}
Yes, as vividly demonstrated in Figure \ref{fig:lr_standard_plot}, the parameter $\lambda$ affects the two metrics in completely different ways. 
The minimum DCF (blue line) remains remarkably flat across the entire parameter sweep, hovering consistently around $0.364$. There are no significant differences in terms of the model's potential discriminatory power as regularization increases. 

Conversely, there is a severe deterioration in terms of calibration (red line). The model maintains a reasonable, albeit imperfect, calibration for low values of $\lambda \le 10^{-3}$. However, as regularization increases beyond $10^{-2}$, the actDCF skyrockets rapidly, eventually maxing out at a worst-case value of $1.0000$ for $\lambda \ge 1$. This clearly shows that $\lambda$ has a much stronger, destructive impact on the reliability of the scores (actDCF) than on their relative ranking ability (minDCF).

\subsection*{How does the regularization coefficient affect the two metrics?}
The regularization coefficient $\lambda$ affects the two metrics in fundamentally different ways due to its mathematical interaction with the Logistic Regression objective function:
\begin{itemize}
    \item \textbf{Effect on minDCF:} The impact is negligible. The minimum DCF evaluates the ability of the classifier to correctly rank samples from highest to lowest probability of belonging to the target class. Because our training dataset has plenty of samples relative to the small number of features (6 dimensions), the model is not prone to overfitting in this linear configuration. Therefore, imposing an $L_2$ penalty does not change the optimal orientation of the separating hyperplane. The relative ordering of the samples is preserved, which keeps the minDCF flat.
    
    \item \textbf{Effect on actDCF:} Regularization has a severe negative impact on actDCF. Theoretically, Logistic Regression minimizes the cross-entropy to produce outputs that can be mapped to Log-Likelihood Ratios. However, the regularization term $\frac{\lambda}{2} ||w||^2$ penalizes the magnitude of the weight vector $w$. As $\lambda$ grows large ($\lambda \to \infty$), the optimization forces the weights to shrink towards zero ($w \to 0$). 
    
    When the weights are extremely small, the raw scores $s(x) = w^T x + b$ become dominated by the bias term $b$, losing their variance and pulling the computed LLRs toward zero. For a highly unbalanced target application like $\pi_T = 0.1$, the optimal theoretical Bayes threshold is high ($t = -\log \frac{0.1}{0.9} \approx 2.19$). If the L2 penalty artificially shrinks all test scores well below this threshold, the model will constantly predict the negative class for \textit{every single sample}, leading to a False Negative Rate of $100\%$ and an actDCF of $1.0000$. This completely destroys the calibration.
\end{itemize}